\hypertarget{fonctions-productives}{%
\section{Fonctions productives}\label{fonctions-productives}}

\hypertarget{valeurs-de-retour}{%
\subsection{Valeurs de retour}\label{valeurs-de-retour}}

Certaines des fonctions intégrées que nous avons utilisées, telles que
les fonctions mathématiques, produisent des résultats.
L\textquotesingle appel de la fonction génère une valeur, que nous
affectons généralement à une variable ou que nous utilisons dans une
expression.

\begin{Shaded}
\begin{Highlighting}[]
\NormalTok{e }\OperatorTok{=}\NormalTok{ math.exp(}\FloatTok{1.0}\NormalTok{)}
\NormalTok{hauteur }\OperatorTok{=}\NormalTok{ rayon }\OperatorTok{*}\NormalTok{ math.sin(radians)}
\end{Highlighting}
\end{Shaded}

Mais jusqu\textquotesingle à présent, aucune des fonctions que nous
avons écrites n\textquotesingle a retourné de valeur.

Dans ce chapitre, nous allons écrire des fonctions productives
(\emph{fruitful functions}). Le premier exemple est \texttt{aire}, qui
retourne l\textquotesingle aire d\textquotesingle un cercle avec le
rayon donné :

\begin{Shaded}
\begin{Highlighting}[]
\ImportTok{import}\NormalTok{ math}

\KeywordTok{def}\NormalTok{ aire(rayon):}
\NormalTok{    a }\OperatorTok{=}\NormalTok{ math.pi }\OperatorTok{*}\NormalTok{ rayon}\OperatorTok{**}\DecValTok{2}
    \ControlFlowTok{return}\NormalTok{ a}
\end{Highlighting}
\end{Shaded}

Nous avons vu l\textquotesingle instruction \texttt{return} auparavant,
mais dans une fonction productive, l\textquotesingle instruction
\texttt{return} inclut une expression. Cette instruction signifie : «
Retourne immédiatement de cette fonction et utilise
l\textquotesingle expression suivante comme valeur de retour. »
L\textquotesingle expression fournie peut être arbitrairement
compliquée, nous pourrions donc écrire cette fonction de manière plus
concise :

\begin{Shaded}
\begin{Highlighting}[]
\KeywordTok{def}\NormalTok{ aire(rayon):}
    \ControlFlowTok{return}\NormalTok{ math.pi }\OperatorTok{*}\NormalTok{ rayon}\OperatorTok{**}\DecValTok{2}
\end{Highlighting}
\end{Shaded}

D\textquotesingle un autre côté, les variables temporaires comme
\texttt{a} peuvent rendre le débogage plus facile.

Parfois, il est utile d\textquotesingle avoir plusieurs instructions de
retour, une dans chaque branche d\textquotesingle une condition :

\begin{Shaded}
\begin{Highlighting}[]
\KeywordTok{def}\NormalTok{ valeur\_absolue(x):}
    \ControlFlowTok{if}\NormalTok{ x }\OperatorTok{\textless{}} \DecValTok{0}\NormalTok{:}
        \ControlFlowTok{return} \OperatorTok{{-}}\NormalTok{x}
    \ControlFlowTok{else}\NormalTok{:}
        \ControlFlowTok{return}\NormalTok{ x}
\end{Highlighting}
\end{Shaded}

Puisque ces instructions \texttt{return} sont dans une condition
alternative, une seule d\textquotesingle entre elles sera exécutée.

Dès qu\textquotesingle une instruction de retour
s\textquotesingle exécute, la fonction se termine sans exécuter les
instructions suivantes. Tout code qui apparaît après une instruction
\texttt{return}, ou à tout autre endroit que le flux
d\textquotesingle exécution ne peut jamais atteindre, est appelé du code
mort (\emph{dead code}).

Dans une fonction productive, c\textquotesingle est une bonne idée de
s\textquotesingle assurer que chaque chemin possible à travers le
programme atteint une instruction \texttt{return}. Par exemple :

\begin{Shaded}
\begin{Highlighting}[]
\KeywordTok{def}\NormalTok{ valeur\_absolue(x):}
    \ControlFlowTok{if}\NormalTok{ x }\OperatorTok{\textless{}} \DecValTok{0}\NormalTok{:}
        \ControlFlowTok{return} \OperatorTok{{-}}\NormalTok{x}
    \ControlFlowTok{if}\NormalTok{ x }\OperatorTok{\textgreater{}} \DecValTok{0}\NormalTok{:}
        \ControlFlowTok{return}\NormalTok{ x}
\end{Highlighting}
\end{Shaded}

Cette fonction est incorrecte car si \texttt{x} se trouve être 0, aucune
des conditions n\textquotesingle est vraie, et la fonction se termine
sans toucher à une instruction \texttt{return}. Si le flux
d\textquotesingle exécution atteint la fin d\textquotesingle une
fonction, la valeur de retour est \texttt{None}, ce qui
n\textquotesingle est pas la valeur absolue de 0.

\begin{Shaded}
\begin{Highlighting}[]
\OperatorTok{\textgreater{}\textgreater{}\textgreater{}} \BuiltInTok{print}\NormalTok{(valeur\_absolue(}\DecValTok{0}\NormalTok{))}
\VariableTok{None}
\end{Highlighting}
\end{Shaded}

\hypertarget{duxe9veloppement-incruxe9mental}{%
\subsection{Développement
incrémental}\label{duxe9veloppement-incruxe9mental}}

Au fur et à mesure que vous écrirez des fonctions plus grandes, vous
constaterez peut-être que vous passez plus de temps à déboguer. Pour
faire face à des programmes de plus en plus complexes, vous voudrez
peut-être essayer un processus appelé développement incrémental
(\emph{incremental development}).

L\textquotesingle objectif du développement incrémental est
d\textquotesingle éviter les longues sessions de débogage en ajoutant et
en testant seulement une petite quantité de code à la fois.

À titre d\textquotesingle exemple, supposons que vous souhaitiez trouver
la distance entre deux points, donnés par leurs coordonnées
(x\textsubscript{1}, y\textsubscript{1}) et (x\textsubscript{2},
y\textsubscript{2}). Par le théorème de Pythagore, la distance est :

\[distance = \sqrt{(x_2 - x_1)^2 + (y_2) - Y_1)^2}\]

La première étape consiste à considérer à quoi devrait ressembler la
fonction \texttt{distance} en Python. Autrement dit, quelles sont les
entrées (paramètres) et quelle est la sortie (valeur de retour) ?

Dans ce cas, les entrées sont deux points, que vous pouvez représenter à
l\textquotesingle aide de quatre nombres. La valeur de retour est la
distance, qui est une valeur à virgule flottante.

Vous pouvez déjà écrire le squelette de la fonction :

\begin{Shaded}
\begin{Highlighting}[]
\KeywordTok{def}\NormalTok{ distance(x1, y1, x2, y2):}
    \ControlFlowTok{return} \FloatTok{0.0}
\end{Highlighting}
\end{Shaded}

Évidemment, cette version ne calcule pas les distances ; elle retourne
toujours zéro. Mais elle est syntaxiquement correcte, et elle
s\textquotesingle exécutera, ce qui signifie que vous pouvez la tester
avant de la rendre plus compliquée.

Pour tester la nouvelle fonction, appelez-la avec des exemples
d\textquotesingle arguments :

\begin{Shaded}
\begin{Highlighting}[]
\OperatorTok{\textgreater{}\textgreater{}\textgreater{}}\NormalTok{ distance(}\DecValTok{1}\NormalTok{, }\DecValTok{2}\NormalTok{, }\DecValTok{4}\NormalTok{, }\DecValTok{6}\NormalTok{)}
\FloatTok{0.0}
\end{Highlighting}
\end{Shaded}

J\textquotesingle ai choisi ces valeurs pour que la distance horizontale
soit 3 et la distance verticale soit 4 ; de cette façon, le résultat est
5 (l\textquotesingle hypoténuse d\textquotesingle un triangle 3-4-5).
Lors du test d\textquotesingle une fonction, il est utile de connaître
la bonne réponse.

À ce stade, nous avons confirmé que la fonction est syntaxiquement
correcte, et nous pouvons commencer à ajouter du code. Après chaque
changement, nous testons à nouveau le programme. Si une erreur survient
à n\textquotesingle importe quel point, nous savons où elle se trouve :
elle doit se trouver dans la dernière ligne que nous avons ajoutée.

Une étape logique de calcul consiste à trouver les différences
x\textsubscript{2} - x\textsubscript{1} et y\textsubscript{2} -
y\textsubscript{1}. Nous stockerons ces valeurs dans des variables
temporaires et les imprimerons :

\begin{Shaded}
\begin{Highlighting}[]
\KeywordTok{def}\NormalTok{ distance(x1, y1, x2, y2):}
\NormalTok{    dx }\OperatorTok{=}\NormalTok{ x2 }\OperatorTok{{-}}\NormalTok{ x1}
\NormalTok{    dy }\OperatorTok{=}\NormalTok{ y2 }\OperatorTok{{-}}\NormalTok{ y1}
    \BuiltInTok{print}\NormalTok{(}\StringTok{\textquotesingle{}dx est\textquotesingle{}}\NormalTok{, dx)}
    \BuiltInTok{print}\NormalTok{(}\StringTok{\textquotesingle{}dy est\textquotesingle{}}\NormalTok{, dy)}
    \ControlFlowTok{return} \FloatTok{0.0}
\end{Highlighting}
\end{Shaded}

Si la fonction fonctionne, elle devrait afficher \texttt{dx\ est\ 3} et
\texttt{dy\ est\ 4}. Si c\textquotesingle est le cas, nous savons que la
fonction obtient les bons arguments et effectue le premier calcul
correctement. Si ce n\textquotesingle est pas le cas, il
n\textquotesingle y a que quelques lignes à vérifier.

Ensuite, nous calculons la somme des carrés de \texttt{dx} et
\texttt{dy} :

\begin{Shaded}
\begin{Highlighting}[]
\KeywordTok{def}\NormalTok{ distance(x1, y1, x2, y2):}
\NormalTok{    dx }\OperatorTok{=}\NormalTok{ x2 }\OperatorTok{{-}}\NormalTok{ x1}
\NormalTok{    dy }\OperatorTok{=}\NormalTok{ y2 }\OperatorTok{{-}}\NormalTok{ y1}
\NormalTok{    carre\_dx }\OperatorTok{=}\NormalTok{ dx}\OperatorTok{**}\DecValTok{2}
\NormalTok{    carre\_dy }\OperatorTok{=}\NormalTok{ dy}\OperatorTok{**}\DecValTok{2}
    \BuiltInTok{print}\NormalTok{(}\StringTok{\textquotesingle{}carre\_dx est\textquotesingle{}}\NormalTok{, carre\_dx)}
    \BuiltInTok{print}\NormalTok{(}\StringTok{\textquotesingle{}carre\_dy est\textquotesingle{}}\NormalTok{, carre\_dy)}
    \ControlFlowTok{return} \FloatTok{0.0}
\end{Highlighting}
\end{Shaded}

Encore une fois, vous devez exécuter le programme à ce stade et vérifier
la sortie (qui devrait être 25).

Enfin, vous pouvez utiliser la fonction \texttt{math.sqrt} pour calculer
et retourner le résultat :

\begin{Shaded}
\begin{Highlighting}[]
\ImportTok{import}\NormalTok{ math}

\KeywordTok{def}\NormalTok{ distance(x1, y1, x2, y2):}
\NormalTok{    dx }\OperatorTok{=}\NormalTok{ x2 }\OperatorTok{{-}}\NormalTok{ x1}
\NormalTok{    dy }\OperatorTok{=}\NormalTok{ y2 }\OperatorTok{{-}}\NormalTok{ y1}
\NormalTok{    carre\_dx }\OperatorTok{=}\NormalTok{ dx}\OperatorTok{**}\DecValTok{2}
\NormalTok{    carre\_dy }\OperatorTok{=}\NormalTok{ dy}\OperatorTok{**}\DecValTok{2}
\NormalTok{    resultat }\OperatorTok{=}\NormalTok{ math.sqrt(carre\_dx }\OperatorTok{+}\NormalTok{ carre\_dy)}
    \ControlFlowTok{return}\NormalTok{ resultat}
\end{Highlighting}
\end{Shaded}

Si cela fonctionne correctement, vous avez terminé. Sinon, vous voudrez
peut-être afficher la valeur de \texttt{resultat} avant
l\textquotesingle instruction de retour.

La version finale de la fonction n\textquotesingle affiche rien
lorsqu\textquotesingle elle est exécutée ; elle retourne simplement une
valeur. Les instructions \texttt{print} que nous avons écrites sont
utiles pour le débogage, mais une fois que vous avez vérifié que la
fonction fonctionne, vous devez les supprimer. Un code comme celui-ci
est appelé code d\textquotesingle échafaudage (\emph{scaffolding}) car
il est utile pour construire le programme mais ne fait pas partie du
produit final.

\hypertarget{composition}{%
\subsection{Composition}\label{composition}}

Comme vous devriez vous y attendre à ce jour, vous pouvez appeler une
fonction à partir d\textquotesingle une autre. Cette capacité est
appelée composition (\emph{composition}).

À titre d\textquotesingle exemple, nous allons écrire une fonction qui
prend deux points, le centre du cercle et un point sur le périmètre, et
qui calcule l\textquotesingle aire du cercle.

Supposons que le point central est stocké dans les variables \texttt{xc}
et \texttt{yc}, et le point de périmètre est dans \texttt{xp} et
\texttt{yp}. La première étape consiste à trouver le rayon du cercle,
qui est la distance entre les deux points. Nous venons
d\textquotesingle écrire une fonction, \texttt{distance}, qui fait cela
:

\begin{Shaded}
\begin{Highlighting}[]
\NormalTok{rayon }\OperatorTok{=}\NormalTok{ distance(xc, yc, xp, yp)}
\end{Highlighting}
\end{Shaded}

La deuxième étape consiste à trouver l\textquotesingle aire
d\textquotesingle un cercle avec ce rayon; nous venons
d\textquotesingle écrire cela aussi :

\begin{Shaded}
\begin{Highlighting}[]
\NormalTok{resultat }\OperatorTok{=}\NormalTok{ aire(rayon)}
\end{Highlighting}
\end{Shaded}

Encapsuler ces étapes dans une fonction produit ceci :

\begin{Shaded}
\begin{Highlighting}[]
\KeywordTok{def}\NormalTok{ aire\_du\_cercle(xc, yc, xp, yp):}
\NormalTok{    rayon }\OperatorTok{=}\NormalTok{ distance(xc, yc, xp, yp)}
\NormalTok{    resultat }\OperatorTok{=}\NormalTok{ aire(rayon)}
    \ControlFlowTok{return}\NormalTok{ resultat}
\end{Highlighting}
\end{Shaded}

Les variables temporaires \texttt{rayon} et \texttt{resultat} sont
utiles pour le développement et le débogage, mais une fois que le
programme fonctionne, nous pouvons le rendre plus concis en composant
les appels de fonctions :

\begin{Shaded}
\begin{Highlighting}[]
\KeywordTok{def}\NormalTok{ aire\_du\_cercle(xc, yc, xp, yp):}
    \ControlFlowTok{return}\NormalTok{ aire(distance(xc, yc, xp, yp))}
\end{Highlighting}
\end{Shaded}

\hypertarget{fonctions-booluxe9ennes}{%
\subsection{Fonctions booléennes}\label{fonctions-booluxe9ennes}}

Les fonctions peuvent retourner des booléens, ce qui est souvent
pratique pour cacher des tests complexes à l\textquotesingle intérieur
de fonctions. Par exemple :

\begin{Shaded}
\begin{Highlighting}[]
\KeywordTok{def}\NormalTok{ est\_divisible(x, y):}
    \ControlFlowTok{if}\NormalTok{ x }\OperatorTok{\%}\NormalTok{ y }\OperatorTok{==} \DecValTok{0}\NormalTok{:}
        \ControlFlowTok{return} \VariableTok{True}
    \ControlFlowTok{else}\NormalTok{:}
        \ControlFlowTok{return} \VariableTok{False}
\end{Highlighting}
\end{Shaded}

Il est courant de donner aux fonctions booléennes des noms qui
ressemblent à des questions par oui ou par non; \texttt{est\_divisible}
retourne \texttt{True} ou \texttt{False} pour indiquer si \texttt{x} est
divisible par \texttt{y}.

Voici un exemple :

\begin{Shaded}
\begin{Highlighting}[]
\OperatorTok{\textgreater{}\textgreater{}\textgreater{}}\NormalTok{ est\_divisible(}\DecValTok{6}\NormalTok{, }\DecValTok{4}\NormalTok{)}
\VariableTok{False}
\OperatorTok{\textgreater{}\textgreater{}\textgreater{}}\NormalTok{ est\_divisible(}\DecValTok{6}\NormalTok{, }\DecValTok{3}\NormalTok{)}
\VariableTok{True}
\end{Highlighting}
\end{Shaded}

Le résultat de l\textquotesingle opérateur \texttt{==} est un booléen,
nous pouvons donc écrire la fonction de manière plus concise en le
retournant directement :

\begin{Shaded}
\begin{Highlighting}[]
\KeywordTok{def}\NormalTok{ est\_divisible(x, y):}
    \ControlFlowTok{return}\NormalTok{ x }\OperatorTok{\%}\NormalTok{ y }\OperatorTok{==} \DecValTok{0}
\end{Highlighting}
\end{Shaded}

Les fonctions booléennes sont souvent utilisées dans des instructions
conditionnelles :

\begin{Shaded}
\begin{Highlighting}[]
\ControlFlowTok{if}\NormalTok{ est\_divisible(x, y):}
    \BuiltInTok{print}\NormalTok{(}\StringTok{\textquotesingle{}x est divisible par y\textquotesingle{}}\NormalTok{)}
\end{Highlighting}
\end{Shaded}

Il peut être tentant d\textquotesingle écrire quelque chose comme :

\begin{Shaded}
\begin{Highlighting}[]
\ControlFlowTok{if}\NormalTok{ est\_divisible(x, y) }\OperatorTok{==} \VariableTok{True}\NormalTok{:}
    \BuiltInTok{print}\NormalTok{(}\StringTok{\textquotesingle{}x est divisible par y\textquotesingle{}}\NormalTok{)}
\end{Highlighting}
\end{Shaded}

Mais la comparaison supplémentaire est inutile.

\hypertarget{glossaire}{%
\subsection{Glossaire}\label{glossaire}}

variable temporaire temporary variable Une variable utilisée pour
stocker une valeur intermédiaire dans un calcul complexe.

code mort dead code Une partie d\textquotesingle un programme qui ne
peut jamais être exécutée, souvent parce qu\textquotesingle elle
apparaît après une instruction de retour.

développement incrémental incremental development Un plan de
développement de programme destiné à éviter le débogage en ajoutant et
en testant seulement une petite quantité de code à la fois.

code d\textquotesingle échafaudage scaffolding Code qui est utilisé
pendant le développement du programme mais qui ne fait pas partie du
produit final.

fonction booléenne boolean function Une fonction qui retourne un
booléen.

\hypertarget{exercices}{%
\subsection{Exercices}\label{exercices}}

\textbf{Exercice 1}

Écrivez une fonction \texttt{compare} qui prend deux valeurs, \texttt{x}
et \texttt{y}, et qui retourne \texttt{1} si
\texttt{x\ \textgreater{}\ y}, \texttt{0} si \texttt{x\ ==\ y}, et
\texttt{-1} si \texttt{x\ \textless{}\ y}.

\textbf{Exercice 2}

Utilisez un développement incrémental pour écrire une fonction nommée
\texttt{hypotenuse} qui retourne la longueur de
l\textquotesingle hypoténuse d\textquotesingle un triangle rectangle,
étant donné la longueur des deux autres côtés. Enregistrez chaque étape
du processus de développement au fur et à mesure que vous avancez.

\textbf{Exercice 3}

Écrivez une fonction booléenne \texttt{est\_entre(x,\ y,\ z)} qui
retourne \texttt{True} si \texttt{x\ ≤\ y\ ≤\ z} et \texttt{False}
sinon.
